\documentclass[10pt,a4paper]{amsart}
\usepackage{amsmath,amssymb,amsthm,graphicx}
% Retained minimal preamble for the Stacks Project chapter more-algebra.tex (More on Algebra).
% The chapter file pinned-original-source.tex carries only \input{preamble} on its first line, so no
% preamble text lives inside the chapter itself.  The Stacks Project uses one single shared preamble for
% every chapter; the snapshot stages it identically as sources/stacks-chow/preamble.tex and
% sources/stacks-etale-cohomology/preamble.tex.  The block below is copied mechanically, line for line,
% from sources/stacks-chow/preamble.tex lines 14-18 (the Xy-pic setup) and 154-217 (the native theorem
% declarations together with the mathematical macros).  Nothing was added, removed or redefined; the
% wrapper class and the amsmath/amssymb/amsthm/graphicx base packages are supplied by the assembler.
% Chapter: more-algebra.tex -- More on Algebra.
% For commutative diagrams we use Xy-pic
\usepackage[all]{xy}

% We use 2cell for 2-commutative diagrams.
\xyoption{2cell}

% Theorem environments.
%
\theoremstyle{plain}
\newtheorem{theorem}[subsection]{Theorem}
\newtheorem{proposition}[subsection]{Proposition}
\newtheorem{lemma}[subsection]{Lemma}

\theoremstyle{definition}
\newtheorem{definition}[subsection]{Definition}
\newtheorem{example}[subsection]{Example}
\newtheorem{exercise}[subsection]{Exercise}
\newtheorem{situation}[subsection]{Situation}

\theoremstyle{remark}
\newtheorem{remark}[subsection]{Remark}
\newtheorem{remarks}[subsection]{Remarks}

\numberwithin{equation}{subsection}

% Macros
%
\def\lim{\mathop{\mathrm{lim}}\nolimits}
\def\colim{\mathop{\mathrm{colim}}\nolimits}
\def\Spec{\mathop{\mathrm{Spec}}}
\def\Hom{\mathop{\mathrm{Hom}}\nolimits}
\def\Ext{\mathop{\mathrm{Ext}}\nolimits}
\def\SheafHom{\mathop{\mathcal{H}\!\mathit{om}}\nolimits}
\def\SheafExt{\mathop{\mathcal{E}\!\mathit{xt}}\nolimits}
\def\Sch{\mathit{Sch}}
\def\Mor{\mathop{\mathrm{Mor}}\nolimits}
\def\Ob{\mathop{\mathrm{Ob}}\nolimits}
\def\Sh{\mathop{\mathit{Sh}}\nolimits}
\def\NL{\mathop{N\!L}\nolimits}
\def\CH{\mathop{\mathrm{CH}}\nolimits}
\def\proetale{{pro\text{-}\acute{e}tale}}
\def\etale{{\acute{e}tale}}
\def\QCoh{\mathit{QCoh}}
\def\Ker{\mathop{\mathrm{Ker}}}
\def\Im{\mathop{\mathrm{Im}}}
\def\Coker{\mathop{\mathrm{Coker}}}
\def\Coim{\mathop{\mathrm{Coim}}}

% Boxtimes
%
\DeclareMathSymbol{\boxtimes}{\mathbin}{AMSa}{"02}

%
% Macros for moduli stacks/spaces
%
\def\QCohstack{\mathcal{QC}\!\mathit{oh}}
\def\Cohstack{\mathcal{C}\!\mathit{oh}}
\def\Spacesstack{\mathcal{S}\!\mathit{paces}}
\def\Quotfunctor{\mathrm{Quot}}
\def\Hilbfunctor{\mathrm{Hilb}}
\def\Curvesstack{\mathcal{C}\!\mathit{urves}}
\def\Polarizedstack{\mathcal{P}\!\mathit{olarized}}
\def\Complexesstack{\mathcal{C}\!\mathit{omplexes}}
% \Pic is the operator that assigns to X its picard group, usage \Pic(X)
% \Picardstack_{X/B} denotes the Picard stack of X over B
% \Picardfunctor_{X/B} denotes the Picard functor of X over B
\def\Pic{\mathop{\mathrm{Pic}}\nolimits}
\def\Picardstack{\mathcal{P}\!\mathit{ic}}
\def\Picardfunctor{\mathrm{Pic}}
\def\Deformationcategory{\mathcal{D}\!\mathit{ef}}

% --- appended 2026-10-12 (Chief): the project own definitions of the reference/slogan/history
% display environments, copied mechanically from the shared Stacks preamble lines 8-12
% (\usepackage{verbatim} plus the three \newenvironment lines, which the chapter preamble keeps
% verbatim from the project preamble). Nothing added, removed or redefined: these environments
% carry no mathematical content and the project itself typesets them as comments. Without these
% three lines a counted statement whose own text uses them fails with Environment slogan undefined.
\usepackage{verbatim}
\newenvironment{reference}{\comment}{\endcomment}
\newenvironment{slogan}{\comment}{\endcomment}
\newenvironment{history}{\comment}{\endcomment}
\usepackage{cleveref}
\crefname{originallemma}{original lemma}{original lemmas}
\Crefname{originallemma}{Original lemma}{Original lemmas}
\makeatletter
\newlabel{lemma-neighbourhood-equivalence}{{24301}{source}{}{sourcepointer-1}{}}
\newlabel{lemma-neighbourhood-equivalence@cref}{{[originallemma][0][]24301}{[1][0][]source}{}{}{}}
\makeatother
\crefname{originalproposition}{original proposition}{original propositions}
\Crefname{originalproposition}{Original proposition}{Original propositions}
\makeatletter
\newlabel{proposition-equivalence}{{24813}{source}{}{sourcepointer-2}{}}
\newlabel{proposition-equivalence@cref}{{[originalproposition][0][]24813}{[1][0][]source}{}{}{}}
\makeatother
\crefname{originalremark}{original remark}{original remarks}
\Crefname{originalremark}{Original remark}{Original remarks}
\makeatletter
\newlabel{remark-glueing-data}{{24641}{source}{}{sourcepointer-3}{}}
\newlabel{remark-glueing-data@cref}{{[originalremark][0][]24641}{[1][0][]source}{}{}{}}
\makeatother
\title{Maps between modules over a flat covering that are I-power torsion in kernel and cokernel come from maps over the base ring}
\author{The Stacks Project Authors}
\date{stacks-project:more-algebra.tex:local-snapshot-20261012; GNU FDL 1.2 or later}
\begin{document}
\maketitle
\setcounter{subsection}{1}
\noindent\textit{Editorial note.} All delivered mathematical text is the author's own, copied from the pinned Stacks Project chapter file more-algebra.tex; no statement or proof step was rewritten, reordered or removed.\par\medskip
\noindent\textit{Editorial note.} The chapter file pinned-original-source.tex carries only \textbackslash{}input\{preamble\} on its first line, so no preamble text lives inside the chapter. The retained file sources/stacks-more-algebra/retained-preamble.tex is a mechanical, line-for-line copy of the project's single shared preamble as staged in the snapshot; nothing was added, removed or redefined and no package was added.\par\medskip
\noindent\textit{Editorial note.} The source URL is the source-level Stacks Project URL, as required when more than one card is produced from one source; the per-claim tag is carried by claim\_key.\par\medskip
\noindent\textit{Editorial note.} The author uses the symbol varphi both for the ambient flat map R -> S and for the given S-module map in the statement and proof; that reuse is reproduced unchanged. The seven references inside the counted proof span are recorded as non-load-bearing navigation pointers.\par\medskip
\noindent\textit{Allowed standard prerequisites (editorial).}\begin{itemize}
\item The gluing category Glue(R -> S, f\_1, ..., f\_t) and the equivalence of categories with R-modules. invoked at proof lines 24904-24910 and 24920-24921. Is the formalism that converts the descended object into an R-module and the morphism into a map over R.
\item I-power torsion modules and their vanishing after localizing at the generators of I. the hypothesis on kernel and cokernel at statement lines 24887-24888 and used at proof line 24903. Makes each localized map an isomorphism, which is what lets the gluing data be assembled.
\item The neighbourhood equivalence identifying the kernel and cokernel of a descended map. invoked at proof lines 24951-24952. Gives the final statement comparing kernel and cokernel of varphi and psi.
\end{itemize}\medskip
\ifdefined\hypertarget\hypertarget{sourcepointer-1}{}\fi
\noindent\textit{Source reference.} Original lemma 24301: sources/stacks-more-algebra/pinned-original-source.tex, source lines 24301--24301. Same-chapter navigation pointer; the statement it names is standard for the section and the argument that consumes it stays inside the counted proof span.\par\medskip
\ifdefined\hypertarget\hypertarget{sourcepointer-2}{}\fi
\noindent\textit{Source reference.} Original proposition 24813: sources/stacks-more-algebra/pinned-original-source.tex, source lines 24813--24813. Same-chapter navigation pointer; the statement it names is standard for the section and the argument that consumes it stays inside the counted proof span.\par\medskip
\ifdefined\hypertarget\hypertarget{sourcepointer-3}{}\fi
\noindent\textit{Source reference.} Original remark 24641: sources/stacks-more-algebra/pinned-original-source.tex, source lines 24641--24641. Same-chapter navigation pointer; the statement it names is standard for the section and the argument that consumes it stays inside the counted proof span.\par\medskip
\begin{lemma}
\label{lemma-application-formal-glueing}
Let $\varphi : R \to S$ be a flat ring map and let $I \subset R$ be a
finitely generated ideal such that $R/I \to S/IS$ is an isomorphism. 
\begin{enumerate}
\item Given an $R$-module $N$, an $S$-module $M'$ and an $S$-module
map $\varphi : M' \to N \otimes_R S$ whose kernel and cokernel are
$I$-power torsion, there exists an $R$-module map
$\psi : M \to N$ and an isomorphism $M \otimes_R S = M'$
compatible with $\varphi$ and $\psi$.
\item Given an $R$-module $M$, an $S$-module $N'$ and an $S$-module
map $\varphi : M \otimes_R S \to N'$ whose kernel and cokernel are
$I$-power torsion, there exists an $R$-module map
$\psi : M \to N$ and an isomorphism $N \otimes_R S = N'$
compatible with $\varphi$ and $\psi$.
\end{enumerate}
In both cases we have $\Ker(\varphi) \cong \Ker(\psi)$ and
$\Coker(\varphi) \cong \Coker(\psi)$.
\end{lemma}

\begin{proof}
Proof of (1). Say $I = (f_1, \ldots, f_t)$. It is clear that
the localization $\varphi_{f_i}$ is an isomorphism.
Thus we see that $(M', N_{f_i}, \varphi_{f_i}, can_{ij})$ is an
object of $\text{Glue}(R \to S, f_1, \ldots, f_t)$, see
Remark \ref{remark-glueing-data}.
By Proposition \ref{proposition-equivalence}
we conclude that there exists an $R$-module $M$ such that
$M' = M \otimes_R S$ and $N_{f_i} = M_{f_i}$ compatibly
with the isomorphisms $\varphi_{f_i}$ and $can_{ij}$. There is a
morphism
$$
(M \otimes_R S, M_{f_i}, can_i, can_{ij}) =
(M', N_{f_i}, \varphi_{f_i}, can_{ij})
\to
(N \otimes_R S, N_{f_i}, can_i, can_{ij})
$$
of $\text{Glue}(R \to S, f_1, \ldots, f_t)$
which uses $\varphi$ in the first component. This
corresponds to an $R$-module map $\psi : M \to N$ (by the equivalence of
categories of Proposition \ref{proposition-equivalence}).
The composition of the base change of $M \to N$ with the
isomorphism $M' \cong M \otimes_R S$ is $\varphi$, in other words
$M \to N$ is compatible with $\varphi$.

\medskip\noindent
Proof of (2). This is just the dual of the argument above.
Namely, the localization $\varphi_{f_i}$ is an isomorphism.
Thus we see that $(N', M_{f_i}, \varphi_{f_i}^{-1}, can_{ij})$ is an
object of $\text{Glue}(R \to S, f_1, \ldots, f_t)$, see
Remark \ref{remark-glueing-data}.
By Proposition \ref{proposition-equivalence}
we conclude that there exists an $R$-module $N$ such that
$N' = N \otimes_R S$ and $N_{f_i} = M_{f_i}$ compatibly
with the isomorphisms $\varphi_{f_i}^{-1}$ and $can_{ij}$. There is a
morphism
$$
(M \otimes_R S, M_{f_i}, can_i, can_{ij}) \to
(N', M_{f_i}, \varphi_{f_i}, can_{ij}) =
(N \otimes_R S, N_{f_i}, can_i, can_{ij})
$$
of $\text{Glue}(R \to S, f_1, \ldots, f_t)$
which uses $\varphi$ in the first component. This
corresponds to an $R$-module map $\psi : M \to N$ (by the equivalence of
categories of Proposition \ref{proposition-equivalence}).
The composition of the base change of $M \to N$ with the
isomorphism $N' \cong N \otimes_R S$ is $\varphi$, in other words
$M \to N$ is compatible with $\varphi$.

\medskip\noindent
The final statement follows for example from
Lemma \ref{lemma-neighbourhood-equivalence}.
\end{proof}

\end{document}
